\section*{अध्याय \ref{ch:b1:weathering-clays} --- अपक्षय और मृत्तिकाएँ}

\begin{solution}{exo:b1:weathering-clays:1}
\ce{K+}: $1/1.38 = 0.72$; \ce{Fe^2+}: $2/0.78 = 2.56$; दोनों 3 से नीचे, घुलनशील धनायन। \ce{Fe^3+}: $3/0.645 = 4.65$; \ce{Ti^4+}: $4/0.605 = 6.61$; दोनों 3 और 10 के बीच, जल-अपघटन-उत्पाद। ऑक्सीजन-रहित जलमग्न मिट्टी में लोहा \ce{Fe^2+} के रूप में मुक्त होता है, जो विलयन में रहता है और जल द्वारा बहा दिया जाता है; अच्छी तरह निकास वाली मिट्टी में ऑक्सीजन उसे \ce{Fe^3+} में बदलती है, जो तुरंत ऑक्साइडों और हाइड्रॉक्साइडों के रूप में अवक्षेपित होकर रहता है।
\end{solution}

\begin{solution}{exo:b1:weathering-clays:2}
A ग्लाइकॉल के साथ फूलता है और तापने पर सिकुड़ता है: एक \omterm{def:b1:weathering-clays:clay-mineral}{स्मेक्टाइट}। B \SI{7}{\angstrom} पर है और \SI{550}{\celsius} पर नष्ट होता है: \omterm{def:b1:weathering-clays:clay-mineral}{कैओलिनाइट}। C, उपचार चाहे जो हो, \SI{10}{\angstrom} पर: \omterm{def:b1:weathering-clays:clay-mineral}{इलाइट}। D, बिना फूले या सिकुड़े \SI{14}{\angstrom} पर: \omterm{def:b1:weathering-clays:clay-mineral}{क्लोराइट}।
\end{solution}

\begin{solution}{exo:b1:weathering-clays:3}
दो ऑर्थोक्लेज़ एक \omterm{def:b1:weathering-clays:clay-mineral}{कैओलिनाइट} के 2 ऐल्युमिनियम देते हैं; $6 - 2 = 4$ सिलिकॉन \ce{H4SiO4} के रूप में निकलता है; 2 \ce{K+} निकलते हैं, 2 \ce{H+} से संतुलित; ऑक्सीजन, $16 + w = 9 + 16$, $w = 9$ देता है:
\[
  \ce{2 KAlSi3O8 + 2 H+ + 9 H2O -> Al2Si2O5(OH)4 + 2 K+ + 4 H4SiO4}.
\]
हाइड्रोजन जाँचिए: $2 + 18 = 20 = 4 + 16$। सारा पोटैशियम निकलता है, और सिलिकॉन का $4/6 =
2/3$ भाग।
\end{solution}

\begin{solution}{exo:b1:weathering-clays:4}
(a) द्रव्यमान नहीं बदलता: $10 \times 999.87/916.72 = \qty{10.91}{cm^3}$।
(b) लगभग \qty{-4}{\celsius} और \qty{-15}{\celsius} के बीच। शून्य से थोड़ा नीचे बर्फ़ दरार को पच्चर से खोलने के लिए पर्याप्त दाब नहीं बना सकती; कठोर पाले में पानी छिद्रों से होकर बर्फ़ की ओर नहीं बहता, और बर्फ़ का बढ़ना रुक जाता है। (c) \cref{prop:b1:weathering-clays:surface} से $n = \qty{1}{m}/\qty{1}{cm} = 100$ के साथ, क्षेत्रफल 100 गुना होता है: \qty{6}{m^2} से \qty{600}{m^2}।
\end{solution}

\begin{solution}{exo:b1:weathering-clays:5}
$S/m = 6/(\rho d)$: \qty{1}{cm} के लिए, $6/(2600\times0.01) =
\qty{0.231}{m^2.kg^{-1}}$; \qty{1}{mm} के लिए, \qty{2.31}{m^2.kg^{-1}}; \qty{1}{\micro m} के लिए,
\qty{2308}{m^2.kg^{-1}}, या \qty{2.3}{m^2.g^{-1}}। रासायनिक अभिक्रियाएँ कणों की सतह पर होती हैं, इसलिए प्रति किलोग्राम चट्टान उनकी दर विशिष्ट सतह के साथ बढ़ती है, सेंटीमीटर से माइक्रोमीटर तक दस हज़ार गुना। \omterm{def:b1:weathering-clays:weathering}{भौतिक अपक्षय} चट्टान को और छोटे टुकड़ों में काटता है और दरारें खोलता है जिनसे पानी प्रवेश करता है: वह उस सतह को गुणित करता है जिस पर \omterm{def:b1:weathering-clays:weathering}{रासायनिक अपक्षय} कार्य कर सकता है।
\end{solution}

\begin{solution}{exo:b1:weathering-clays:6}
एक टन ऐल्बाइट $\num{e6}/262.22 = \qty{3813.6}{mol}$ है। \cref{prop:b1:weathering-clays:feldspars} से, दो मोल एक \omterm{def:b1:weathering-clays:clay-mineral}{कैओलिनाइट} देते हैं:
$\qty{1906.8}{mol} \times \qty{258.16}{g.mol^{-1}} = \qty{492.3}{kg}$। तीन सिलिकॉन परमाणुओं में से दो निकलते हैं: $2\times3813.6 = \qty{7627.2}{mol}$
\ce{SiO2}, \qty{458.2}{kg}। सारा सोडियम निकलता है: $3813.6\times22.99 =
\qty{87.7}{kg}$। फ़ेल्डस्पार का आधे से कम द्रव्यमान मृत्तिका के रूप में बचा रहता है।
\end{solution}

\begin{solution}{exo:b1:weathering-clays:7}
\cref{prop:b1:weathering-clays:immobile} से, एक किलोग्राम ग्रेनाइट छोड़ता है
$0.30/0.60 = \qty{0.5}{kg}$ सैप्रोलाइट। खोया भाग: सोडियम के लिए $1 -
(0.35/3.5)(0.5) = 0.95$; सिलिकॉन $1 - (60/72)(0.5) = 0.583$; ऐल्युमिनियम $1 -
(27/14)(0.5) = 0.036$। ऐल्युमिनियम व्यावहारिक रूप से रुका है, सोडियम लगभग पूरा चला गया है, और सिलिकॉन का लगभग \qty{60}{\percent} भी; पर बचे का बड़ा भाग क्वार्ट्ज में रहता है: फ़ेल्डस्पार \omterm{def:b1:weathering-clays:clay-mineral}{कैओलिनाइट} में बदल रहे हैं, \omterm{def:b1:weathering-clays:hydrolysis-degree}{एक-सिऐलिटीकरण}।
\end{solution}

\begin{solution}{exo:b1:weathering-clays:8}
पाइरोफ़िलाइट: $x = y = 0$, आवेश 0। मस्कोवाइट परत: $x = 1$, आवेश $-1$।
बेइडेलाइट परत: $x = 0.5$, आवेश $-0.5$। मॉन्टमॉरिलोनाइट परत: $y = 0.33$,
आवेश $-0.33$। दो \omterm{def:b1:weathering-clays:clay-mineral}{स्मेक्टाइट} परतें, बेइडेलाइट और मॉन्टमॉरिलोनाइट, फूल सकती हैं: उनका छोटा आवेश कम अंतर-परत धनायनों से संतुलित है, जो परतों को ढीला पकड़ते हैं और जलयोजित होते समय जल खींचते हैं। पाइरोफ़िलाइट के पास जलयोजित करने के लिए कोई अंतर-परत धनायन नहीं है, और मस्कोवाइट का आवेश $-1$ पोटैशियम आयनों से संतुलित है जो परतों को जकड़ते हैं।
\end{solution}

\begin{solution}{exo:b1:weathering-clays:9}
(a) Mg: 2 और 2; Si: 1 और 1; C: 4 और 4; O: $4 + 8 + 4 = 16$ और $12 + 4 =
16$; H: 8 और $4 + 4 = 8$; आवेश: 0 और $4 - 4 = 0$। (b) फ़ॉर्स्टेराइट का मोलर द्रव्यमान
$2\times24.305 + 28.085 + 4\times15.999 =
\qty{140.69}{g.mol^{-1}}$ है, इसलिए एक टन \qty{7108}{mol} है। वह $4\times7108 = \qty{28431}{mol}$ \ce{CO2} उपभुक्त करता है, \qty{1251}{kg}; हर सूत्र-एकक में दो मैग्नीशियम के साथ, \cref{prop:b1:weathering-clays:co2} $2\times7108$ मोल संचित करता है, \qty{626}{kg}, उपभुक्त मात्रा का आधा। (c) एक टन कैल्साइट, \qty{9991}{mol}, घुलते समय \qty{9991}{mol} \ce{CO2} (\qty{440}{kg}) उपभुक्त करता है, और अवक्षेपित होते समय वह सारी लौटा देता है: कुछ संचित नहीं होता।
\end{solution}

\begin{solution}{exo:b1:weathering-clays:10}
(a) $3 < 3/r < 10$ से $\qty{0.3}{\angstrom} < r < \qty{1.0}{\angstrom}$।
(b) $3/0.27 = 11.1 > 10$: बोरॉन ऑक्सीऋण के रूप में चलता है, बोरेट। (c) $2/0.83 =
2.41$, एक घुलनशील धनायन; $4/0.53 = 7.55$, एक जल-अपघटन-उत्पाद। ऑक्सीजन-दरिद्र जल में मैंगनीज़ \ce{Mn^2+} के रूप में घुला होता है; जहाँ स्रोत हवा से मिलता है, वह \ce{Mn^4+} तक ऑक्सीकृत होता है, जो तुरंत चट्टान पर काले ऑक्साइड के रूप में अवक्षेपित होता है।
(d) उसका \omterm{def:b1:weathering-clays:ionic-potential}{आयनिक विभव}, $4/0.40 = 10.0$, जल-अपघटन-उत्पादों और ऑक्सीऋणों के बीच की सीमा पर है: सिलिकॉन न ऐल्युमिनियम की तरह रुकता है न गंधक की तरह बह जाता है, और निक्षालित भाग इस पर निर्भर है कि कितना पानी गुज़रता है, जो \omterm{def:b1:weathering-clays:hydrolysis}{जल-अपघटन} की तीन मात्राओं का आधार है।
\end{solution}

\begin{solution}{exo:b1:weathering-clays:11}
(a) कार्बन के \qty{12.011}{g.mol^{-1}} के साथ, \num{11.7e12} से
\qty{13.6e12}{mol} प्रति वर्ष 0.141 से 0.163 अरब टन कार्बन है; \qty{44.009}{g.mol^{-1}} के साथ, 0.515 से 0.599 अरब टन कार्बन डाइऑक्साइड।
(b) $4.08/13.6 = \qty{30}{\percent}$ से $4.08/11.7 = \qty{35}{\percent}$ तक।
(c) उपभुक्त कार्बन डाइऑक्साइड का आधा संचित होता है जब कैल्शियम और मैग्नीशियम कार्बोनेट के रूप में अवक्षेपित होते हैं: प्रति वर्ष 0.070 से 0.082 अरब टन कार्बन। वास्तविक आँकड़ा छोटा है, क्योंकि बाइकार्बोनेट का एक भाग सोडियम और पोटैशियम के साथ जाता है, जो समुद्र-तल पर कोई कार्बोनेट नहीं बनाते।
\end{solution}

\begin{solution}{exo:b1:weathering-clays:12}
Si/Al $= (w_{\ce{SiO2}}/60.08)/(2w_{\ce{Al2O3}}/101.96)$। A: अब $0.7748/0.7748 =
1.00$, \omterm{def:b1:weathering-clays:clay-mineral}{कैओलिनाइट} का अनुपात, \omterm{def:b1:weathering-clays:hydrolysis-degree}{एक-सिऐलिटीकरण}। B के लिए $0.9987/0.4904 = 2.04$, एक 2:1 \omterm{def:b1:weathering-clays:clay-mineral}{मृत्तिका-खनिज}, \omterm{def:b1:weathering-clays:hydrolysis-degree}{द्वि-सिऐलिटीकरण}। C: $0.0333/1.2358 = 0.027$, लगभग कोई सिलिकॉन नहीं बचा, गिब्साइट, \omterm{def:b1:weathering-clays:hydrolysis-degree}{ऐलिटीकरण}। निक्षालन की सबसे कम से सबसे तीव्र ओर, यानी सबसे गर्म और सबसे आर्द्र जलवायु की ओर: B, A, C।
\end{solution}

\begin{solution}{pb:b1:weathering-clays:1}
\textbf{1.} \ce{2 NaAlSi3O8 + 2 H+ + 9 H2O -> Al2Si2O5(OH)4 + 2 Na+ + 4 H4SiO4}।

\textbf{2.} \ce{CaAl2Si2O8 + 2 H+ + H2O -> Al2Si2O5(OH)4 + Ca^2+}। कोई सिलिकॉन नहीं निकलता: एनॉर्थाइट में पहले से ही \omterm{def:b1:weathering-clays:clay-mineral}{कैओलिनाइट} का प्रति ऐल्युमिनियम एक सिलिकॉन है।

\textbf{3.} एक गिब्साइट एक ऐल्युमिनियम लेता है, एक एनॉर्थाइट का आधा: दोगुना करने पर,
\ce{CaAl2Si2O8 + 2 H+ + 6 H2O -> 2 Al(OH)3 + Ca^2+ + 2 H4SiO4}। ऑक्सीजन:
$8 + 6 = 6 + 8$; हाइड्रोजन: $2 + 12 = 6 + 8$।

\textbf{4.} \ce{5 NaAlSi3O8 + 4 H+ + 16 H2O -> 2 Na_{0.5}Al2(Si_{3.5}Al_{0.5})O10(OH)2 + 4 Na+ + 8 H4SiO4}; % ce-unbalanced-ok: decimal subscripts the checker cannot read; balance tested in tests/bachelor-1/test_weathering_clays.py
हाइड्रोजन, बाईं ओर $4 + 32 = 36$ और दाईं ओर $2\times2 + 8\times4 = 36$, संतुलित है।

\textbf{5.} हर \ce{H+} एक \ce{HCO3-} के साथ आता है, जो एक \ce{CO2} से और एक
\ce{H2O}:
\ce{2 NaAlSi3O8 + 2 CO2 + 11 H2O -> Al2Si2O5(OH)4 + 2 Na+ + 2 HCO3- + 4 H4SiO4}
और इसी तरह \ce{CaAl2Si2O8 + 2 CO2 + 3 H2O -> Al2Si2O5(OH)4 + Ca^2+ + 2 HCO3-}।

\textbf{6.} ऐल्बाइट से बेइडेलाइट बाइसिअलिटीकरण है; ऐल्बाइट और एनॉर्थाइट से
\omterm{def:b1:weathering-clays:clay-mineral}{कैओलिनाइट} \omterm{def:b1:weathering-clays:hydrolysis-degree}{एक-सिऐलिटीकरण} हैं; एनॉर्थाइट से गिब्साइट \omterm{def:b1:weathering-clays:hydrolysis-degree}{ऐलिटीकरण} है।

\textbf{7.} $\num{e6}/262.22 = \qty{3813.6}{mol}$ ऐल्बाइट का;
$\num{e6}/278.20 = \qty{3594.5}{mol}$ एनॉर्थाइट का।

\textbf{8.} ऐल्बाइट: $3813.6/2 = \qty{1906.8}{mol}$ \omterm{def:b1:weathering-clays:clay-mineral}{कैओलिनाइट},
\qty{492.3}{kg}। एनॉर्थाइट: \qty{3594.5}{mol}, \qty{928.0}{kg}।

\textbf{9.} ऐल्बाइट: $2\times3813.6 = \qty{7627.2}{mol}$ of \ce{SiO2},
\qty{458.2}{kg}, और \qty{3813.6}{mol} सोडियम, \qty{87.7}{kg}। एनॉर्थाइट:
कोई सिलिका नहीं, और \qty{3594.5}{mol} कैल्शियम, \qty{144.1}{kg}।

\textbf{10.} ऐल्बाइट: $1/2.62 = \qty{0.382}{m^3}$; उसका \omterm{def:b1:weathering-clays:clay-mineral}{कैओलिनाइट},
$0.4923/2.6 = \qty{0.189}{m^3}$: ठोस अपने आयतन का \qty{50.4}{\percent} खोता है। एनॉर्थाइट: $1/2.75 = \qty{0.364}{m^3}$; उसका \omterm{def:b1:weathering-clays:clay-mineral}{कैओलिनाइट},
$0.9280/2.6 = \qty{0.357}{m^3}$: केवल \qty{1.8}{\percent} की हानि।

\textbf{11.} खाली जगह छिद्र बन जाती है, जो जल और हवा से भरी है: ऐल्बाइट से समृद्ध एक ग्रेनाइट पर एक सैप्रोलाइट हल्की, छिद्रमय और भुरभुरी है, उसके क्वार्ट्ज-कण थोड़ी-सी मृत्तिका में ढीले। एनॉर्थाइट लगभग अपने आयतन के बराबर \omterm{def:b1:weathering-clays:clay-mineral}{कैओलिनाइट} में बदलता है, इसलिए अपक्षीकृत गैब्रो अपने अधिकांश ठोस द्रव्य को रखता है।

\textbf{12.} सोडियम ($z/r = 0.98$) और कैल्शियम (2.0) घुलनशील धनायन हैं और निकलते हैं; ऐल्युमिनियम (5.6) एक जल-अपघटन-उत्पाद है और \omterm{def:b1:weathering-clays:clay-mineral}{कैओलिनाइट} में पूरा रुकता है; सिलिकॉन (10.0), सीमा पर, ऐल्बाइट से आंशिक रूप से निकलता है और एनॉर्थाइट से बिल्कुल नहीं, जिसका सिलिकॉन \omterm{def:b1:weathering-clays:clay-mineral}{कैओलिनाइट} रख सकता है।

\textbf{13.} ऐल्बाइट पर एक \ce{CO2}: \qty{3813.6}{mol}, $3813.6 \times 44.01 =
\qty{167.8}{kg}$।

\textbf{14.} एनॉर्थाइट पर दो: \qty{7189.0}{mol}, \qty{316.4}{kg}।

\textbf{15.} $7189.0 \times 61.02 = \qty{438.7}{kg}$ बाइकार्बोनेट।

\textbf{16.} जीव इसे कैल्शियम कार्बोनेट के रूप में अवक्षेपित करते हैं,
\ce{Ca^2+ + 2 HCO3- -> CaCO3 + CO2 + H2O}: हर कैल्शियम के लिए, दो में से एक
\ce{CO2} हवा में लौटती है, \qty{3594.5}{mol} या \qty{158.2}{kg}, और
दूसरी कार्बोनेट में दबी रहती है।

\textbf{17.} समुद्र-तल पर सोडियम कोई कार्बोनेट नहीं बनाता: उसका बाइकार्बोनेट कार्बोनेट चट्टान के रूप में संचित नहीं हो सकता, और दीर्घकालिक संचयन केवल कैल्शियम और मैग्नीशियम से गिना जाता है।

\textbf{18.} एक टन कैल्साइट घुलते समय \qty{9991}{mol} \ce{CO2} उपभुक्त करता है और अवक्षेपित होते समय वह सारी लौटा देता है: कुछ संचित नहीं होता, और केवल कैल्शियम और मैग्नीशियम सिलिकेटों का अपक्षय ही कार्बन डाइऑक्साइड को सदा के लिए नीचे खींचता है।

\textbf{19.} $\qty{200}{m^3.s^{-1}} \times \qty{0.3}{mol.m^{-3}} =
\qty{60}{mol.s^{-1}}$; \qty{3.156e7}{s} के एक वर्ष में, \qty{1.893e9}{mol}, या
$\num{1.893e9}/\num{50000} = \qty{37870}{mol}$ प्रति वर्ग किलोमीटर।

\textbf{20.} प्रति एनॉर्थाइट दो बाइकार्बोनेट: प्रति वर्ष \qty{9.467e8}{mol},
$\num{9.467e8}\times\qty{278.20}{g} = \qty{263400}{t}$ प्रति वर्ष, या
\qty{5.27}{t} प्रति वर्ग किलोमीटर।

\textbf{21.} खपी ताज़ी चट्टान $5.27/0.20 = \qty{26.3}{t}$ है, या
$\num{26340}/2700 = \qty{9.75}{m^3}$, प्रति वर्ग किलोमीटर और वर्ष: प्रति वर्ष
\qty{9.75}{\micro m} मोटी परत, प्रति दस लाख वर्ष \qty{9.8}{m}।

\textbf{22.} प्रति कैल्शियम एक \ce{CO2} संचित: $\num{9.467e8} \times 44.01 =
\qty{4.17e10}{g}$, लगभग प्रति वर्ष \qty{41700}{t}।

\textbf{23.} यह केवल विलयन में बहाए गए द्रव्य को गिनता है। अपक्षय से बचे \omterm{def:b1:weathering-clays:clay-mineral}{कैओलिनाइट} और क्वार्ट्ज भी अपरदन द्वारा हटाए जाते हैं, नदी में कीचड़ और बालू के रूप में, और वे नहीं गिने जाते; न ही चट्टान के अन्य खनिजों का अपक्षय। और यदि बाइकार्बोनेट का एक भाग सोडियम या पोटैशियम फ़ेल्डस्पारों से आता है, जो प्रति सूत्र-एकक दो के बजाय एक बाइकार्बोनेट देते हैं, तो उसकी आपूर्ति के लिए और चट्टान अपक्षयित होनी पड़ी।

\textbf{24.} \cref{prop:b1:weathering-clays:co2} से, प्रति मोल एनॉर्थाइट एक मोल \ce{CO2} संचित होता है:
$3594.5 \times \qty{44.01}{g} = \qty{158.19}{kg}$। एक टन एनॉर्थाइट का अपक्षय सदा के लिए
\textbf{158}~kg कार्बन डाइऑक्साइड नीचे खींचता है।
\end{solution}
